\subsection{Siegel's Hydra and numen formalism: typed source boundary}
\label{sec:Siegel-Hydra-numen}

The controlling witness in this subsection is Siegel's version-pinned source
for arXiv:2601.17030v3 \cite{Siegel2026v3}.  Its Hydra and numen sections are
unchanged from v2 to v3; the v3 change is confined to the later Fourier
calculation audited at the end of this subsection.  The source permits both
number fields and finite extensions of $\mathbb F_q(T)$.  Accordingly,
$\mathcal O_K$ is a finite free module over the integer ring
$\mathcal O_{\mathbb F}$ of the base field; it has a finite
$\mathbb Z$-basis only in the number-field case.

Let $K/\mathbb F$ be such a global field, let
$\Lambda\subsetneq\mathcal O_K$ be a nonzero ideal, and retain an enumeration
of the $p=\lvert\mathcal O_K/\Lambda\rvert$ cosets by
$\{0,\ldots,p-1\}$, with $0$ denoting $\Lambda$.  If $\ell$ is a place, an
$\mathbb F$-linear map used on $K_\ell$ will always be assumed bounded at
$\ell$.  This assumption is necessary: an arbitrary
$\mathbb F$-linear endomorphism of $K$ need not extend continuously to a fixed
completion.  The operator norm is
\[
 \lVert r\rVert_\ell
 =\sup_{z\in K\setminus\{0\}}\frac{|r(z)|_\ell}{|z|_\ell},
\]
when this supremum is finite.  The source display includes $z=0$ in a quotient
by $|z|_\ell$ and its following remark asserts extendibility without this
boundedness condition.

\begin{definition}[Siegel's printed branch data]
\label{def:Siegel-Hydra-v3}
A pre-Hydra has, for every $j\in\mathcal O_K/\Lambda$, an invertible affine
$\mathbb F$-linear branch
\[
 H_j:K\longrightarrow K,\qquad H_j(z)=r_j(z)+c_j,
 \qquad
 r_j\in\operatorname{End}_{\mathbb F}^{\times}(K),\quad c_j\in K.
\]
The selected piecewise evaluator is
\[
 H(z)=H_{[z]_\Lambda}(z)\qquad(z\in\mathcal O_K).
\]
The source's literal Definition~5.3 additionally requires
$H_j(z)\in\mathcal O_K$ for every $z\in\mathcal O_K$ and every $j$.
\end{definition}

The enumeration in the source should read
$i=\Lambda_{\beta(i)}$ for a coset $i\in\mathcal O_K/\Lambda$, not
$i\in\Lambda_{\beta(i)}$.  After the cosets have been replaced by digit
labels, $\operatorname{DigSum}_p$ uses those labels directly; applying
$\beta$ to them a second time is ill-typed.  We retain the enumeration rather
than normalize it away.

\begin{sourceissue}[the two printed integrality predicates]
\label{issue:Siegel-Hydra-integrality}
Definition~5.3 requires
\[
 H_j(\mathcal O_K)\subseteq\mathcal O_K
 \qquad\text{for every }j,
\tag{G}
\]
whereas Definition~5.4 calls the same Hydra \emph{integral} when
\[
 H_j(z)\in\mathcal O_K
 \quad\Longleftrightarrow\quad z\in j+\Lambda
 \qquad(z\in\mathcal O_K).
\tag{E}
\]
These conditions cannot coexist.  Since $\Lambda$ is proper, choose a coset
different from $j+\Lambda$ and an integer $z$ in it.  Condition~(G) says
$H_j(z)$ is integral and condition~(E) says it is not.  The source's own
shortened Collatz example displays the intended distinction:
$(3z+1)/2$ is integral exactly on odd $z$, so it satisfies~(E) but not~(G).
\end{sourceissue}

We therefore use three explicitly different predicates:
\[
\begin{array}{ll}
\text{\emph{selected-branch admissible}}&
 H_j(j+\Lambda)\subseteq\mathcal O_K,\\[2mm]
\text{\emph{exact-integral on global integers}}&
 H_j^{-1}(\mathcal O_K)\cap\mathcal O_K=j+\Lambda,\\[2mm]
\text{\emph{inverse-integral at a completion $v$}}&
 H_j^{-1}(\mathcal O_{K_v})=j+\Lambda\mathcal O_{K_v}.
\end{array}
\]
The first makes the piecewise evaluator a map
$\mathcal O_K\to\mathcal O_K$.  The second detects the assigned coset among
global integers.  The third is the stronger local inverse-image statement
needed by a backward branch-correctness argument.  It is not implied merely
by writing the second predicate.

\begin{proposition}[scalar Hydra--Matthews bridge at retained modulus]
\label{prop:Siegel-scalar-Hydra-Matthews}
Let $K=\mathbb Q$, $\mathcal O_K=\mathbb Z$, and
$\Lambda=D\mathbb Z$, with the residue labels $0,\ldots,D-1$ retained.
Scalar selected-branch-admissible data
\[
 H_i(n)=a_i n+b_i
\]
are in exact two-sided correspondence with Matthews data
\[
 H_i(n)=\frac{m_i n-s_i}{D},
 \qquad m_i,s_i\in\mathbb Z,\qquad
 s_i\equiv i m_i\pmod D,
\]
by
\[
 m_i=Da_i,\qquad s_i=-Db_i.
\]
The complete integrality locus of the $i$th extension is
\[
 i+\frac{D}{\gcd(m_i,D)}\mathbb Z.
\]
Consequently, exact integrality on global integers is equivalent to
$\gcd(m_i,D)=1$ for every $i$.
\end{proposition}

\begin{proof}
Admissibility gives $H_i(i),H_i(i+D)\in\mathbb Z$.  Subtracting gives
$Da_i\in\mathbb Z$, and then
\[
 Db_i=D H_i(i)-(Da_i)i\in\mathbb Z.
\]
Thus $m_i$ and $s_i$ are integral, while
\[
 m_i i-s_i=D(a_i i+b_i)=D H_i(i)
\]
proves the congruence.  Conversely, that congruence makes
$(m_i n-s_i)/D$ integral on $i+D\mathbb Z$, and the two displayed
assignments are literal inverses.

For arbitrary $n\in\mathbb Z$, branch integrality is equivalent to
$m_i(n-i)\equiv0\pmod D$.  Writing $g_i=\gcd(m_i,D)$, this is equivalent to
$n-i\in(D/g_i)\mathbb Z$, which proves the locus and the final assertion.
\end{proof}

This proposition applies only to scalar branches.  A genuinely non-scalar
$\mathbb F$-linear Hydra is not silently placed in the Matthews or rcwa
classes.  For the scalar subclass, the further reduced rcwa crosswalk is
Proposition~\ref{prop:Matthews-Kohl-crosswalk}; neither step minimizes the
retained modulus.

\begin{sourceissue}[uniform separation is not bounded finiteness]
\label{issue:Siegel-discrete-orbit-classification}
Source Proposition~4.2 claims that if $Y$ is uniformly separated in an
arbitrary valued field, every orbit of every map $T:Y\to Y$ is preperiodic or
norm-divergent.  Its proof uses the stronger assertion that every bounded
intersection with $Y$ is finite.

For a nontrivially valued counterexample, take $K=\mathbb Q((t))$ with the
$t$-adic absolute value, $Y=\mathbb Z\subset K$, and $T(n)=n+1$.  Distinct
integers have distance one, but the orbit of zero is injective and has norm
one after the first term.  It is neither preperiodic nor divergent.

The proof becomes valid under the exact replacement
\[
 \#\{y\in Y:|y|\leq R\}<\infty
 \qquad\text{for every }R>0.
\tag{BF}
\]
For a number field, $\mathcal O_K$ is boundedly finite in its full Minkowski
embedding.  It need not be discrete at one archimedean embedding:
$\mathbb Z[\sqrt2]$ has nonzero elements arbitrarily close to zero there.
Nor does escape in the full Minkowski norm imply escape at every individual
archimedean absolute value.
\end{sourceissue}

For the numen, let $\Sigma_p^*$ be the finite words in
$\{0,\ldots,p-1\}$, including the empty word.  The source fixes
least-significant digit first:
\[
 \operatorname{DigSum}_p(j_1,\ldots,j_N)
 =\sum_{a=1}^{N}j_a p^{a-1}.
\]

\begin{proposition}[finite affine word and numen cocycle]
\label{prop:Siegel-word-numen}
For $\mathbf j=(j_1,\ldots,j_N)$, define
\[
 H_{\mathbf j}=H_{j_1}\circ\cdots\circ H_{j_N},
 \qquad
 H_{\mathbf j}(z)=M_{\mathbf j}(z)+X_{\mathbf j}.
\]
Then
\[
 M_{\mathbf j}=r_{j_1}\circ\cdots\circ r_{j_N},
 \qquad
 X_{\mathbf j}
 =\sum_{m=0}^{N-1}
   (r_{j_1}\circ\cdots\circ r_{j_m})(c_{j_{m+1}}),
\tag{W}
\label{eq:Siegel-word-formula}
\]
with the empty operator product equal to $I$.  For concatenation,
\[
 M_{\mathbf i\wedge\mathbf j}
 =M_{\mathbf i}\circ M_{\mathbf j},
 \qquad
 X_{\mathbf i\wedge\mathbf j}
 =M_{\mathbf i}(X_{\mathbf j})+X_{\mathbf i}.
\tag{C}
\label{eq:Siegel-word-cocycle}
\]
\end{proposition}

\begin{proof}
The rightmost branch in $H_{j_1}\circ\cdots\circ H_{j_N}$ acts first.
Composing two affine maps gives
\[
 (M_i,X_i)\circ(M_j,X_j)
 =(M_i\circ M_j,M_i(X_j)+X_i).
\]
Induction proves both~\eqref{eq:Siegel-word-formula} and
\eqref{eq:Siegel-word-cocycle}; the empty word gives $(I,0)$.
\end{proof}

Here and below a word is an address.  It is not automatically an actual
piecewise-Hydra itinerary.

\begin{proposition}[centered numen descent and multiplier non-descent]
\label{prop:Siegel-centered-numen-descent}
Let $\mathbf i\sim\mathbf j$ when their DigSum values agree.  If
$c_0=H_0(0)=0$, then $X_{\mathbf i}=X_{\mathbf j}$ whenever
$\mathbf i\sim\mathbf j$, so the finite numen factors as
\[
 \Sigma_p^*/{\sim}\ \xrightarrow[\cong]{\operatorname{DigSum}_p}\mathbb N_0
 \xrightarrow{X_H}K.
\]
The multiplier does not generally factor through this quotient:
\[
 M_{\mathbf j\wedge(0)}=M_{\mathbf j}\circ r_0
\]
need not equal $M_{\mathbf j}$.
\end{proposition}

\begin{proof}
Every finite DigSum fibre consists of the canonical base-$p$ word and its
finitely many terminal-zero extensions.  Since the appended zero is the
rightmost map in the composition,
\[
 X_{\mathbf j\wedge(0)}
 =H_{\mathbf j}(H_0(0))=H_{\mathbf j}(0)=X_{\mathbf j}.
\]
Iteration proves descent.  The displayed multiplier identity follows directly
from Proposition~\ref{prop:Siegel-word-numen} and supplies the obstruction.
\end{proof}

Thus every occurrence of $M_H(n)$ requires a canonical digit word and its
retained length.  The source proves descent only for $X_H$.

\begin{proposition}[finite digit functional equations]
\label{prop:Siegel-numen-functional-equations}
Functions $f:\mathbb N_0\to K$ satisfying
\begin{equation}
\label{eq:Siegel-finite-recursion}
 f(pn+j)=r_j(f(n))+c_j
 \qquad(n\in\mathbb N_0,\ 0\leq j<p)
\end{equation}
are in bijection with the affine fibre
\[
 \{a\in K:(I-r_0)a=c_0\},
\]
by $a=f(0)$.  If $I-r_0$ is invertible, there is one solution, with
$a=(I-r_0)^{-1}c_0$.  If
$n=\sum_{m=0}^{N-1}d_m p^m$ is its canonical digit expansion and
\[
 P_0=I,\qquad P_m=r_{d_0}\circ\cdots\circ r_{d_{m-1}},
\]
then
\begin{equation}
\label{eq:Siegel-corrected-finite-truncation}
 f(n)=P_N(a)+\sum_{m=0}^{N-1}P_m(c_{d_m}).
\end{equation}
\end{proposition}

\begin{proof}
Putting $n=j=0$ in~\eqref{eq:Siegel-finite-recursion} gives
$(I-r_0)f(0)=c_0$.  Conversely, once an admissible $a$ is chosen, repeated
Euclidean division by $p$ reaches zero and recursively defines one value at
every nonnegative integer.  Applying the recursion along the canonical digits
proves~\eqref{eq:Siegel-corrected-finite-truncation}.
\end{proof}

The finite theorem is sound.  It also isolates the term omitted in the
source's proof of its continuation lemma.

\begin{sourceissue}[Lemma 6.1 as printed]
\label{issue:Siegel-numen-extension}
Lemma~6.1 assumes only that $I-r_0$ is invertible, but its finite formula uses
the centered word sum and omits $P_N(a)$ from
\eqref{eq:Siegel-corrected-finite-truncation}.  Its equation~(6.29) prints
equality between the norm of a sum and the sum of termwise norm bounds; only
$\leq$ follows.  Its pointwise density definition uses a limsup frequency
when $\log\lVert r_j\rVert_\ell\leq0$ and a liminf when that logarithm is
positive.  These are the reverse of the choices needed for an upper Lyapunov
bound.  Almost-everywhere convergence leaves a null nonconvergence set on
which no value is assigned.  The proof establishes uniqueness on
$\mathbb N_0$, not uniqueness of arbitrary functions on $\mathbb Z_p$
satisfying the same recursion.  Finally, in the nonarchimedean case the
displayed bound omits $|a|_\ell$.

The last omission has a concrete instance.  Over $\mathbb Q_2$, take
$r_0=3$ and $c_0=1$.  Then
$a=(1-3)^{-1}=-1/2$, so $|a|_2=2>|c_0|_2=1$.
\end{sourceissue}

\begin{theorem}[corrected numen continuation]
\label{thm:Siegel-numen-extension-repair}
Fix a completion $K_\ell$ at which all $r_j$ are bounded, assume
$I-r_0$ is invertible, and put $a=(I-r_0)^{-1}c_0$.  For
\[
 \mathfrak z=\sum_{m\geq0}d_m p^m\in\mathbb Z_p,
\]
define
\begin{equation}
\label{eq:Siegel-corrected-continuation}
 f_N(\mathfrak z)
 =P_N(a)+\sum_{m=0}^{N-1}P_m(c_{d_m}),
 \qquad
 P_m=r_{d_0}\circ\cdots\circ r_{d_{m-1}}.
\end{equation}
For $q_{j,m}=m^{-1}\#\{0\leq k<m:d_k=j\}$ and
$\alpha_j=\log\lVert r_j\rVert_\ell$, set
\begin{equation}
\label{eq:Siegel-upper-Lyapunov}
 L^+(\mathfrak z)=
 \sum_{\alpha_j<0}\alpha_j\liminf_{m\to\infty}q_{j,m}
 +\sum_{\alpha_j>0}\alpha_j\limsup_{m\to\infty}q_{j,m}.
\end{equation}
Then:
\begin{enumerate}[label=\textnormal{(\roman*)}]
\item if $L^+(\mathfrak z)<0$, the sequence
  \eqref{eq:Siegel-corrected-continuation} converges;
\item if $\prod_{j=0}^{p-1}\lVert r_j\rVert_\ell<1$, it converges for
  Haar-almost every $\mathfrak z$ and defines a measurable function on its
  conull convergence domain;
\item if $\rho=\max_j\lVert r_j\rVert_\ell<1$, it converges uniformly on
  $\mathbb Z_p$ to a continuous function;
\item whenever $\mathfrak z$ lies in the convergence domain,
  $p\mathfrak z+j$ also lies there and
  \[
   f(p\mathfrak z+j)=r_j(f(\mathfrak z))+c_j;
  \]
\item if $\ell$ is nonarchimedean and
  $\max_j\lVert r_j\rVert_\ell\leq1$, every finite truncation satisfies
  \[
   |f_N(\mathfrak z)|_\ell
   \leq\max\{|a|_\ell,\max_j|c_j|_\ell\}.
  \]
\end{enumerate}
In case~(iii), the continuous extension is the unique continuous function
with the digit recursion and initial value $a$.
\end{theorem}

\begin{proof}
Submultiplicativity gives
\[
 \lVert P_m\rVert_\ell
 \leq\prod_{j=0}^{p-1}
      \lVert r_j\rVert_\ell^{m q_{j,m}}.
\]
For a negative $\alpha_j$, the limsup of
$\alpha_jq_{j,m}$ is $\alpha_j\liminf q_{j,m}$; for a positive
$\alpha_j$, it is $\alpha_j\limsup q_{j,m}$.  Since there are finitely many
digits,
\[
 \limsup_{m\to\infty}\frac1m\log\lVert P_m\rVert_\ell
 \leq L^+(\mathfrak z).
\]
Under~(i), there are $\eta<1$ and $m_0$ such that
$\lVert P_m\rVert_\ell\leq\eta^m$ for $m\geq m_0$.  The translations form a
finite set, so the series in
\eqref{eq:Siegel-corrected-continuation} is absolutely convergent and
$P_N(a)\to0$.

Haar-almost every base-$p$ digit sequence has each digit frequency $1/p$.
Its Lyapunov exponent is bounded above by
\[
 \frac1p\sum_j\log\lVert r_j\rVert_\ell
 =\frac1p\log\prod_j\lVert r_j\rVert_\ell,
\]
which proves~(ii); a pointwise limit of measurable truncations is measurable
on the convergence domain.  Under~(iii),
$\lVert P_m\rVert_\ell\leq\rho^m$ for every digit sequence, and the geometric
tail estimate is uniform.  The truncations are locally constant, hence their
uniform limit is continuous.

The finite identity
\[
 f_N(p\mathfrak z+j)=r_j(f_{N-1}(\mathfrak z))+c_j
\]
proves~(iv) on passage to the limit.  In the nonarchimedean case, the
ultrametric inequality and $\lVert P_m\rVert_\ell\leq1$ give~(v).
Finally, $\mathbb N_0$ is dense in $\mathbb Z_p$, and
Proposition~\ref{prop:Siegel-numen-functional-equations} fixes all its values;
two continuous extensions therefore agree everywhere.
\end{proof}

Functional equations alone do not give literal uniqueness among arbitrary
functions on $\mathbb Z_p$.  For example, with $p=2$,
$r_0=r_1=1/2$, and $c_0=c_1=0$, choose a non-eventually-periodic digit
sequence.  Its tail-equivalence class is countable and contains no shift
cycle.  A nonzero seed value propagates consistently over that class by the
invertible branch factors, while the function may be set to zero off the
class.  The result is Borel and satisfies the recursion, but is not the zero
function.

\begin{proposition}[operator-typed periodic word]
\label{prop:Siegel-periodic-word-fixed-point}
For a nonempty finite word $\mathbf w$, put
$M=M_{\mathbf w}$ and $X=X_{\mathbf w}$.  Then
\begin{equation}
\label{eq:Siegel-repeated-word}
 X_{\mathbf w^{\wedge m}}
 =\sum_{k=0}^{m-1}M^k(X).
\end{equation}
If $\lVert M\rVert_\ell<1$, then $I-M$ is invertible on $K_\ell$ and
\[
 z_{\mathbf w}=(I-M)^{-1}X
\]
is the unique fixed point of the affine extension $H_{\mathbf w}$.
It is a point of an actual $N$-step piecewise-Hydra cycle only if, for
$\mathbf w=(w_1,\ldots,w_N)$,
\begin{equation}
\label{eq:Siegel-branch-correctness}
\begin{split}
 z_{\mathbf w}&\in w_N+\Lambda,\\
 H_{w_N}(z_{\mathbf w})&\in w_{N-1}+\Lambda,\\
 &\ \vdots\\
 H_{w_2}\circ\cdots\circ H_{w_N}(z_{\mathbf w})
   &\in w_1+\Lambda.
\end{split}
\end{equation}
\end{proposition}

\begin{proof}
The cocycle gives
$X_{\mathbf w^{\wedge m}}=M(X_{\mathbf w^{\wedge(m-1)}})+X$;
induction proves~\eqref{eq:Siegel-repeated-word}.  The Neumann series
$\sum_{k\geq0}M^k$ converges in operator norm and is the inverse of $I-M$.
Thus $M(z_{\mathbf w})+X=z_{\mathbf w}$, and subtracting two fixed-point
equations proves uniqueness.  Since the rightmost branch acts first,
conditions~\eqref{eq:Siegel-branch-correctness} are exactly the successive
branch-selection requirements.  Without them, the calculation concerns only
the affine extension.
\end{proof}

The source's fractions $(1-M^m)/(1-M)$ and $X/(1-M)$ must be replaced by
operator composition with $(I-M)^{-1}$.  Its equality
$\lVert M^m\rVert=\lVert M\rVert^m$ is false for a general operator norm;
the sufficient relation is
$\lVert M^m\rVert\leq\lVert M\rVert^m$.

\begin{sourceissue}[Siegel Theorem 6.1 as stated]
\label{thm:Siegel-correspondence-registered}
The source assumes a proper, centered, integral $\Lambda$-Hydra and the two
conditions
\begin{enumerate}
\item for every $n\geq1$, some place $\ell_n$ satisfies
  $\lVert M_H(n)\rVert_{\ell_n}<1$;
\item some nontrivial nonarchimedean place $v$ satisfies
  $\lVert r_j\rVert_v>1$ for every digit $j$ and $v(\Lambda)\geq1$.
\end{enumerate}
It then states:
\begin{enumerate}
\item an element $z\in\mathcal O_K$ is periodic for $H$ if and only if
  $z=X_H(\mathfrak z)$ for some
  $\mathfrak z\in(\mathbb Z_p\setminus\mathbb N_0)\cap\mathbb Q$;
\item if $K$ is a number field and
  $X_H(\mathfrak z)\in\mathcal O_K$ for an irrational
  $\mathfrak z\in\mathbb Z_p\setminus\mathbb Q$, then the source calls
  $X_H(\mathfrak z)$ a trajectory divergent with respect to every
  nontrivial archimedean absolute value.
\end{enumerate}
This is a source-statement record, not an admitted theorem of this edition.
\end{sourceissue}

The proof sketch does not certify either conclusion.  The exact blocking
points are as follows.
\begin{enumerate}
\item The literal Hydra and integral hypotheses have no simultaneous models,
  by Source Issue~\ref{issue:Siegel-Hydra-integrality}.
\item $M_H$ does not descend through finite DigSum equivalence.  Thus
  $M_H(n)$ is not defined until a canonical word and its retained length are
  specified.  The displayed $\lambda_p(n)$ is likewise not defined.
\item Two formulas use an undefined $M_q$.  Scalar fractions such as
  $(1-M^m)/(1-M)$ are not typed for general endomorphisms, and the equality
  $\lVert M^m\rVert=\lVert M\rVert^m$ is false.  The operator-typed,
  submultiplicative replacement is
  Proposition~\ref{prop:Siegel-periodic-word-fixed-point}.
\item Passing from an arbitrary finite word to the canonical digit word of
  its DigSum may remove terminal zeroes.  Repeating those two words need not
  give the same multiplier or affine map.  The centered all-zero word and
  the fixed point zero also lie outside the displayed $n\geq1$ construction.
\item The notation $v(\Lambda)$, the \emph{correctness} predicate, and the
  claimed extension of the branches to $\mathcal O_{K_v}$ are not defined or
  proved in this source.  A formal affine fixed point is not an actual
  piecewise orbit until every condition in
  \eqref{eq:Siegel-branch-correctness} is established.
\item The converse uses integrality backwards: from an integral later value
  it infers that a possibly nonintegral predecessor is integral.  The
  printed global predicate does not imply the completion-level inverse-image
  identity needed for this step, nor is a local-to-global arrow supplied.
\item The forward direction gives, at most, an existential rational address
  for a periodic point.  It does not prove that the same numen fibre contains
  no irrational address.  Consequently it cannot support the later inference
  that an irrational address has a nonpreperiodic value.  The zero fibre is
  also unclassified.
\item Part II begins by invoking a nonexistent hypothesis~(iii), then imports
  a zero-fibre assertion and a correctness result without rebuilding their
  hypotheses.  It also depends on the false uniformly-discrete orbit
  classification repaired in
  Source Issue~\ref{issue:Siegel-discrete-orbit-classification}.
\item A numen value is a point, not a trajectory.  Even after replacing it
  by its forward orbit, escape in the full Minkowski embedding would not by
  itself prove divergence at every individual archimedean place.
\end{enumerate}

Accordingly, the periodic-word fixed-point proposition is retained at its
proved scope, while the source's rational-fibre correspondence and
irrational-address divergence conclusions are not proved and are not used.
The legacy audit identifier \texttt{PO-COL-000006} records that gap; it is not
a mathematical problem or a commitment attributed to the source or reader.

\begin{sourceissue}[the v3 nonarchimedean Fourier inversion]
\label{issue:Siegel-v3-prop76}
Let $L=K_\ell$ be a finite extension of $\mathbb Q_q$, let
$\mathcal O_L$ be its integer ring, let $\pi$ be a uniformizer, and let
$\mathfrak D=\mathfrak D_{L/\mathbb Q_q}$ be its different.  The source uses
the trace character
\[
 e_\ell(x)=
 \exp\!\left(2\pi i
 \{\operatorname{Tr}_{L/\mathbb Q_q}(x)\}_q\right).
\]
The annihilator of $\mathcal O_L$ under the pairing
$(t,x)\mapsto e_\ell(tx)$ is
\[
 \mathfrak D^{-1}
 =\{t\in L:\operatorname{Tr}_{L/\mathbb Q_q}(t\mathcal O_L)
                    \subseteq\mathbb Z_q\},
\]
not $\mathcal O_L$ in general.  Thus the source's asserted dual
$L/\mathcal O_L$ is correct for this character only after a
conductor normalization, for example when the relevant different is
trivial.  With the printed character the dual is
$L/\mathfrak D^{-1}$.

For comparison, if an $L$-valued random variable $X$ has first been proved
to lie in $\mathcal O_L$ almost surely, define
\[
 \widehat\mu(t)=\mathbb E[e_\ell(-tX)]
 \qquad(t\in L/\mathfrak D^{-1}).
\]
Writing $e$ and $f=[L:\mathbb Q_q]/e$ for the ramification index and residue
degree, finite character orthogonality gives the correctly typed identity
\begin{equation}
\label{eq:Siegel-corrected-local-Fourier}
\Pr(X\in w+\pi^n\mathcal O_L)
=q^{-nf}
 \sum_{t\in\pi^{-n}\mathfrak D^{-1}/\mathfrak D^{-1}}
   \widehat\mu(t)e_\ell(tw)
\qquad(w\in\mathcal O_L,\ n\geq0).
\end{equation}
Indeed, the displayed frequency quotient is the complete character group of
$\mathcal O_L/\pi^n\mathcal O_L$, which has $q^{nf}$ elements.  Averaging
those characters is one on the zero residue class and zero on every other
class, proving~\eqref{eq:Siegel-corrected-local-Fourier}.

This correction does not certify source Proposition~7.6.  Its surrounding
Fourier chain has further independent defects:
\begin{enumerate}
\item for a nonscalar bounded local branch $r_j$, the frequency map is the
  trace adjoint $r_j^\dagger$, defined by
  \[
   \operatorname{Tr}_{L/\mathbb Q_q}(t\,r_jx)
   =\operatorname{Tr}_{L/\mathbb Q_q}(r_j^\dagger t\,x),
  \]
  not $t\mapsto r_jt$ unless the required self-adjoint or scalar identity is
  proved;
\item the characteristic function is naturally defined first on $L$; its
  descent to a quotient requires an almost-sure support statement matched to
  the annihilator;
\item source Proposition~7.5 multiplies values in an object assumed only to
  be an abelian group, shadows its summation variable, and leaves another
  variable unbound, so it cannot justify the later rearrangement;
\item $B_{H,\ell}=e\max_j\log_q|c_j|_\ell$ is undefined when every
  translation is zero and may be negative, although Proposition~7.5 is
  invoked only for a nonnegative integer parameter;
\item the correct size relation is
  $q^{B/e}=|\pi^{-B}|_\ell$, not $|\pi^B|_\ell$; a real absolute value cannot
  be compared to the field element $\pi^B$; and division by $\pi^B$ cannot
  be replaced by division by the real number $q^{B/e}$.
\end{enumerate}
The visible v3 changes repair neither the dual lattice nor these typing and
scaling defects.  No proposition depending on this Fourier chain is used
downstream.  The legacy audit identifier
\texttt{PO-COL-}\allowbreak\texttt{000007} records the resulting dependency
gap; it does not convert the defective chain into an open mathematical problem.
\end{sourceissue}

\begin{nonclaim}
The finite numen is an address-to-state construction.  It is not identified
with a state-to-itinerary conjugacy, and its fibres are not classified.
Neither the corrected continuation theorem nor the formal periodic-word
fixed point proves an ordinary-integer Collatz endpoint.  The source's
Correspondence Principle and Fourier consequences are recorded as open
audits, not imported conclusions.
\end{nonclaim}
